% !TEX root = ../notes.tex

\documentclass[../notes.tex]{subfiles}

\begin{document}

\section{April 1}
Here we go.

\subsection{Back to Modular Forms}
We now explain how modular forms give rise to automorphic forms.
\begin{exe}
	Set $G\coloneqq\op{SL}_{2,\QQ}$, fix a level $K_f\subseteq G(\AA_\QQ^\infty)$, and set $\Gamma\coloneqq K_f\cap\op{SL}_2(\QQ)$. Now, we let $M_k(\Gamma)$ be the holomorphic modular forms of weight $k$ and level $\Gamma$, and we define a map $\varphi_\bullet M_k(\Gamma)\to\mc A$ as follows. Identify $\mc A^{K_f}\subseteq C^\infty(G(\QQ)\backslash G(\AA_\QQ)/K_f)$ with $C^\infty(\Gamma\backslash G(\RR))$ by strong approximation. Then given $f\in M_k(\Gamma)$, we define $\varphi_f\in C^\infty(\Gamma\backslash G(\RR))$ by
	\[\varphi_f\left(\begin{bmatrix}
		a & b \\ c & d
	\end{bmatrix}\right)=f\left(\frac{ai+b}{ci+d}\right)(ci+d)^{-k}.\]
	We claim that $\varphi_f$ is an automorphic form.
\end{exe}
\begin{proof}
	By construction, we see that $\varphi_f$ is $K_f$-invariant. Here are the other checks for $\varphi_f$ to be an automorphic form.
	\begin{itemize}
		\item Because $\varphi$ has weight $k$, we see that $\varphi_f(gs)=s^k\cdot\varphi_f(g)$ for any $s\in\op{SO}_2(\RR)$ (identified with $S^1$). Thus, $\varphi_f$ is an eigenfunction under $K_\infty$ of weight $k$, so $K_\infty$-finiteness follows.
		\item Let $L$ be the Cauchy--Riemann operator. Then $f$ being holomorphic means that $Lf=0$. On the other hand, it turns out that $L$ is the action by $f\in\mf{sl}_2(\CC)$, so it follows that $\Omega f=\frac{k(k-2)}2f$, so $Z\mf g$-finiteness follows.
		\item Given the above, $f$ being holomorphic at the cusps implies that $\varphi_f$ has moderate growth by an examination of the $q$-expansion at $\infty$ in $\mc H$.
		\qedhere
	\end{itemize} 
\end{proof}
\begin{remark} \label{rem:mk-gamma-as-automorphic}
	This argument shows that $M_k(\Gamma)$ embeds into the automorphic forms which are killed by the ideal $\left(\Omega-k(k-2)/2\right)\subseteq Z\mf g$, have $K_\infty$-action given by $\sigma\colon\op{SO}_2(\RR)\to S^1$ which is the $k$th power, and they are invariant under $K_f$. This subspace of automorphic forms must in fact be given by $M_k(\Gamma)$: indeed, one can define $f$ by reversing the definition of $\varphi_f$ (which is well-defined due to the weight $k$), and having weight $k$ and having the given action by Casimir implies that $f$ is holomorphic.
\end{remark}
\begin{remark}
	If one changes $\frac{k(k-2)}2$ to some other number $\lambda$, then we get Maass forms. They are notably not holomorphic!
\end{remark}
While we're here, we restate our admissibility theorem.
\begin{notation}
	Fix a semisimple group $G$ over a number field $F$. Further, choose $J\subseteq Z\mf g$ of finite codimension, an irreducible representation $\sigma$ of $K_\infty$, and a level $K_f\subseteq G(\AA_F^\infty)$. Then we define $\mc A(J,\sigma,K_f)$ to be the space of automorphic forms killed by $J$ with $K_\infty$-type $\sigma$, and fixed by $K_f$.
\end{notation}
\begin{theorem}[Harish-Chandra] \label{thm:hc-admissible-automorphic}
	Fix a semisimple group $G$ over a number field $F$. Further, choose $J\subseteq Z\mf g$ of finite codimension, an irreducible representation $\sigma$ of $K_\infty$, and a level $K_f\subseteq G(\AA_F^\infty)$. Then
	\[\dim\mc A(J,\sigma,K_f)<\infty.\]
\end{theorem}
By varying $\sigma$ and $K_f$, one recovers \Cref{thm:automorphic-get-admissible}. Let's work this out for modular forms.
\begin{example} \label{ex:dim-mod-form}
	\Cref{thm:hc-admissible-automorphic} implies that $M_k(\Gamma)$ is finite-dimensional by \Cref{rem:mk-gamma-as-automorphic}. Let's show this directly. Note that $f\in M_k(\Gamma)$ gives rise to a differential form $f(z)\,(dz)^{k/2}\in\omega_{\mc H/\CC}^{\otimes k/2}$. (This makes sense if $k$ is even. If $k$ is odd, then we should view $\omega$ as $\OO(-2)|_{\mc H}$, where $\OO(-2)$ comes from the embedding $\mc H\into\PP^1$; then $\omega^{\otimes k/2}$ is just $\OO(-k)|_{\mc H}$.) The same argument as in \Cref{ex:bounded-at-infinity} with $q$-expansions shows that $f(z)\,(dz)^\ell$ is allowed a pole of order at most $k/2$ at each cusp. Thus, we have embedded $M_k(\Gamma)$ into the global sections of a coherent sheaf (namely, differential forms with bounded poles) on the compact manifold $\Gamma\backslash\mc H^*$. The compact complex curve is a projective algebraic curve, and this vector bundle is also algebraic (it is some power of $\omega$, twisted appropriately), so its space of global sections is finite-dimensional!
\end{example}
\begin{remark}
	The proof of \Cref{ex:dim-mod-form} shows that modular forms are algebraic, so we may find a basis defined over a fixed number field.
\end{remark}

\subsection{Suddenly Shimura Varieties}
It seems difficult to compute the dimension in \Cref{thm:hc-admissible-automorphic}. One approach is to embed $\Gamma\backslash G(\RR)/K_\infty$ into a Zariski open subset of a complex projective variety and then proceed as in \Cref{ex:dim-mod-form} to compute dimensions via algebraic geometry. For example, this is done in the theory of Shimura varieties, where $G(\RR)/K_\infty$ is found to be a Hermitian symmetric domain.
\begin{example}
	For $G=\op{Sp}(V,\omega)$, we have $K_\infty=\op U(V)$; set $X\coloneqq\op{Sp}(V)/K_\infty$. Note that $\dim_\RR X=n(n+1)$, which is even, so this has a chance of admitting a complex structure. In fact, $X$ can be embedded into the space of Lagrangian subspaces in the standard symplectic space $V_\CC$: indeed, $\op{Sp}(V)$ certainly acts on isotropic subspaces of $V_\CC$, and one can check that this action is transitive. To see that we can achieve $\op U(n)$ as a stabilizer, recall that $\op U(n)$ is embedded as automorphisms of $V$, which we give a complex structure $J$, which then gives rise to a unique Hermitian structure $h$ for which $\im h$ is the symplectic form on $V$. Now, $V$ with the complex structure $J$ allows us to decompose $V_\CC=V\oplus\overline V$ (via the decomposition $\CC\otimes\CC=\CC\oplus\CC$). Now, $V\subseteq V_\CC$ is Lagrangian, and its stabilizer is exactly $K_\infty=\op U(V,h)$!
\end{example}
\begin{example}
	However, this approach cannot work in general: $\op{SL}_3(\ZZ)\backslash\op{SL}_3(\RR)/\op{SO}_3(\RR)$ is five-dimen\-sional, so it cannot be turned into a complex projective variety.
\end{example}
\begin{remark}
	For $G$ simple, one finds that $X\coloneqq G(\RR)/K_\infty$ is a Hermitian symmetric domain if and only if $Z(K_\infty)$ contains a copy of $S^1$. Intuitively, what is going on here is that this action of $S^1$ will rotate tangent vectors in $G(\RR)/K_\infty$, thereby giving a complex structure on $T_1X$.
\end{remark}
\begin{example}
	The real form $G(\RR)=\op{SU}(p,q)$ has $K_\infty=\op S(\op U(p)\times\op U(q))$ has a copy of $\op U(1)$, so this case admits a Hermitian symmetric domain. This case turns out to embed into the Grassmannian variety of $p$-dimensional subspaces of $\RR^{p+q}$.
\end{example}

\subsection{Siegel Domains}
For a semisimple group $G$, one can tile the double quotient $G(F)\backslash G(\AA_F)/K_f$ by things which look like $\Gamma\backslash G(F_v)$. It may be desirable to put such elements into a ``normal form.''
\begin{theorem}[Iwasawa decomposition] \label{thm:iwasawa}
	Fix a real split reductive group $G$ over $\RR$. Then
	\[G(\RR)=N(\RR)A(\RR)^\circ K_\infty,\]
	where $N(\RR)$ is the unipotent subgroup of a Borel subgroup, $A(\RR)$ has positive diagonal elements from a maximal split torus, and $K_\infty$ is the maximal compact subgroup. In fact, this decomposition is unique.
\end{theorem}
\begin{proof}
	If $G=\op{SL}_n$, then one does Gram--Schmidt orthogonalization. The proof in general is hard.
\end{proof}
The moral is that $\Gamma\backslash G(F_v)/K_\infty$ is covered by the domain $(\Gamma\cap N(\RR))\backslash N(\RR)A(\RR)^\circ$.
\begin{example}
	For $G=\op{SL}_2(\RR)$, we see that $\Gamma\cap N(\RR)$ consists of translations, and $N(\RR)A(\RR)^\circ$ is the upper-half plane. Thus, the quotient is a vertical strip, which has infinite volume.
\end{example}
Having infinite volume is understood to be undesirable, so we want to cut down our covering set more. The idea is to make $A$ smaller.
\begin{definition}[Siegel domain]
	Fix a split semisimple group $G$ over $\RR$, and choose an arithmetic subgroup $\Gamma\subseteq G(\RR)$. Let $N\subseteq G$ be the unipotent radical of a Borel subgroup, and let $A\subseteq G$ be a split torus. Further, fix a compact fundamental domain $\Omega_N\subseteq N(\RR)$ for the action of $\Gamma\cap N(\RR)$, and choose some $\varepsilon>0$ and define $A_\varepsilon\subseteq A(\RR)^\circ$ by
	\[A_\varepsilon\coloneqq\{a\in A(\RR)^\circ:\alpha_i(a)\ge\varepsilon\},\]
	where $\alpha_i$ is lifted to a character $A(\RR)\to\RR^\times$. Then a \textit{Siegel domain} is the image of $\Omega_NA_\varepsilon$ in $G(\RR)/K_\infty$.
\end{definition}
\begin{remark} \label{rem:siegel-domain-exists}
	It turns out that Siegel domains always exist so that $\Omega_NA_\varepsilon$ covers $\Gamma\backslash G(\RR)/K_\infty$. We show this on the homework for $\op{SL}_n$.
\end{remark}
Here is our application.
\begin{theorem}
	Fix a semisimple group $G$ over a number field $F$. Given $G(F)\backslash G(\AA_F)$ a measure by pushing forward the Haar measure on $G(\AA_F)$. Then the volume of $G(F)\backslash G(\AA_F)$ is finite.
\end{theorem}
\begin{proof}
	We will only prove this in the case where $G$ is split and $F=\QQ$, but the difficulties to remove these assumptions are not important. Because $K_f$ is open and compact, it is enough to show that $G(\QQ)\backslash G(\AA_\QQ)/K_f$ has finite volume, for any level $K_f\subseteq G(\AA_\QQ)$. Recall that $G(\QQ)\backslash G(\AA_\QQ)$ can be partitioned into pieces which look like $\Gamma\backslash G(\RR)$, for possibly varying $\Gamma$.

	Thus, we would like to show that $\Gamma\backslash G(\RR)$ has finite volume. By \Cref{rem:siegel-domain-exists}, it is enough to show that $\Omega_NA_\varepsilon K_\infty$ has finite volume, where $\Omega_NA_\varepsilon$ is a Siegel domain. Now, because $G$ is unimodular, \Cref{thm:iwasawa} allows us to decompose the Haar measure $dg$ into $d_\ell(na)d_rk$ (indeed, this is a measure, and it has enough invariance requirements). Further, $N$ is normal and unimodular in $AN$, so it turns out that
	\[d_\ell(na)=da\cdot dn=a^{-2\rho}\,dn\,da.\]
	Namely, $a^{-2\rho}$ comes out as the modular character from $A$ acting on $N$.
	\begin{example}
		Take $G=\op{SL}_2(\RR)$. With $x$ as the coordinate of $\begin{bsmallmatrix}
			1 & x \\ & 1
		\end{bsmallmatrix}$ in $N(\RR)$, and $y$ as the coordinate of $\op{diag}(\sqrt y,1/\sqrt y)$ in $A(\RR)^\circ$, one finds that the Haar measure on $G$ is
		\[\frac{dx\,dy}{y^2}\,dk,\]
		where $dk$ is the Haar measure on $K_\infty$.
	\end{example}
	Thus, the volume of $\Omega_NA_\varepsilon K_\infty$ is the product of the volumes of $\Omega_N$, $A_\varepsilon$, and $K_\infty$, all considered separately, along with some factor of $a^{-2\rho}$. Notably, the volumes of $\Omega_N$ and $K_\infty$ are finite because these are compact, so it only remains to show that
	\[\int_{A_\varepsilon}a^{-2\rho}\,da\]
	is finite. Now, $A$ is split, so this turns into a product of integrals over $\mathbb G_m$s, so we reduce to the case of $\op{SL}_2(\RR)$, where the integral is
	\[\int_\varepsilon^\infty a^{-1}\,\frac{da}a<\infty,\]
	so we are done!
\end{proof}
\begin{example}
	One can actually compute that
	\[\op{vol}(\op{SL}_n(\ZZ)\backslash\op{SL}_n(\RR))=\zeta(2)\zeta(3)\cdots\zeta(n),\]
	provided that the left-hand volume is normalized appropriately (given by taking a top form of $\mf{sl}_n(\RR)$).
\end{example}

\end{document}